
\paragraph{\textbf{3.5.}}
\SGARefTarget{sga3:xx:3.5}
Before stating the following theorem, let us make a few remarks.  As
usual, we identify \(\mathfrak g^{-\alpha}\) with
\((\mathfrak g^\alpha)^{\otimes-1}\).  Similarly, we identify
\[
  \operatorname{Hom}_{\mathscr O_S}
    (\mathfrak g^{-\alpha},\mathfrak g^\alpha)
  \quad\text{with}\quad
  (\mathfrak g^\alpha)^{\otimes2},
\]
and hence
\[
  \operatorname{Isom}_{\mathscr O_S\text{-mod.}}
    \bigl(W(\mathfrak g^{-\alpha}),W(\mathfrak g^\alpha)\bigr)
  \simeq
  W\bigl((\mathfrak g^\alpha)^{\otimes2}\bigr)^\times.
\]
If \(w\in N^\times(S)\), then
\(\operatorname{Ad}(w)\) permutes
\(\mathfrak g^\alpha\) and \(\mathfrak g^{-\alpha}\)
\textup{(\SGARefLink{sga3:xx:3.1}{3.1} \textup{\SGARefLink{sga3:xx:3.1:item:iii}{(iii)}})}, and therefore defines an isomorphism
\[
  a_\alpha(w)\colon
  \mathfrak g^{-\alpha}\xrightarrow{\ \sim\ }\mathfrak g^\alpha.
\]
We thus identify it with a section
\[
  a_\alpha(w)\in
  \Gamma\bigl(S,(\mathfrak g^\alpha)^{\otimes2}\bigr)^\times.
\]
This construction is compatible with base change and therefore
defines a morphism
\[
  a_\alpha\colon
  N^\times\longrightarrow
  W\bigl((\mathfrak g^\alpha)^{\otimes2}\bigr)^\times
\]
such that
\[
  a_\alpha(w)Y=\operatorname{Ad}(w)Y
\]
for every \(w\in N^\times(S')\),
\(Y\in\Gamma(S',\mathfrak g^{-\alpha})^\times\), and \(S'\to S\).

\paragraph{\textbf{Theorem 3.6.}}
\SGARefTarget{sga3:xx:3.6}
\emph{\textup{(i)} For every \(S'\to S\),
\(w\in N^\times(S')\), and
\(Y\in W(\mathfrak g^{-\alpha})(S')\), one has}
\[
  \intaut(w)\exp(Y)=\exp\bigl(a_\alpha(w)Y\bigr).
\]

\smallskip
\noindent\emph{\textup{(ii)} One has}
\[
  a_\alpha(tw)=\alpha(t)a_\alpha(w),
  \qquad
  a_\alpha(wt)=\alpha(t)^{-1}a_\alpha(w).
\]

\smallskip
\noindent\emph{\textup{\SGAInlineRefTarget{sga3:xx:3.6:item:iii}(iii)} If one similarly defines}
\[
  a_{-\alpha}\colon
  N^\times\longrightarrow
  W\bigl((\mathfrak g^{-\alpha})^{\otimes2}\bigr)^\times,
\]
\emph{then}
\[
  a_{-\alpha}(w)=a_\alpha(w)^{-1}.
\]
\footnote{That is, \(a_{-\alpha}(w)\) and \(a_\alpha(w)\) are paired;
cf.\ \SGARefLink{sga3:xx:2.6.1}{2.6.1}.}

\smallskip
\noindent\emph{\textup{\SGAInlineRefTarget{sga3:xx:3.6:item:iv}(iv)} For every
\(X\in W(\mathfrak g^\alpha)^\times(S')\), \(S'\to S\), one has}
\[
  a_\alpha\bigl(w_\alpha(X)\bigr)=-X^2.
\]

\smallskip
\noindent\emph{Proof.}
Assertion \textup{\SGARefLink{sga3:xx:1.5:item:i}{(i)}} is immediate from the characterization of the
morphisms \(\exp\) given in \SGARefLink{sga3:xx:1.5}{1.5}.  Assertion \textup{\SGARefLink{sga3:xx:1.5:item:ii}{(ii)}} is
immediate, as is \textup{\SGARefLink{sga3:xx:3.6:item:iii}{(iii)}}.  Let us prove \textup{\SGARefLink{sga3:xx:3.6:item:iv}{(iv)}}.  Let
\[
  X\in\Gamma(S',\mathfrak g^\alpha)^\times,
  \qquad
  Z\in\Gamma(S',\mathfrak g^\alpha).
\]
By definition,\footnote{The editors supplied the details that follow.}
\[
\begin{aligned}
  a_\alpha\bigl(w_\alpha(X)\bigr)^{-1}(Z)
    &={}\operatorname{Ad}\bigl(w_\alpha(X)\bigr)(Z)\\
    &={}\operatorname{Ad}(\exp(X))
       \operatorname{Ad}(\exp(-X^{-1}))
       \operatorname{Ad}(\exp(X))(Z).
\end{aligned}
\]
Applying formulas \textup{(2\(^{\prime}\))} and \textup{(2)} of
\SGARefLink{sga3:xx:2.10}{Lemma 2.10}, together with
\[
  H_{-\alpha}=-H_\alpha,
  \qquad
  \bar\alpha(H_\alpha)=2
\]
\textup{(loc.\ cit., \textup{\SGARefLink{sga3:xx:lemma:2-10:equation:4}{(4)}} and \textup{\SGARefLink{sga3:xx:lemma:2-10:equation:5}{(5)}})}, and
\(\langle X,X^{-1}\rangle=1\) \textup{(\SGARefLink{sga3:xx:2.6}{2.6})}, the expression on the
right is, successively,
\[
\begin{aligned}
  &\operatorname{Ad}(\exp(X))
    \operatorname{Ad}(\exp(-X^{-1}))(Z)\\
  &\quad={}\operatorname{Ad}(\exp(X))
       \bigl(Z+\langle X^{-1},Z\rangle(H_\alpha-X^{-1})\bigr)\\
  &\quad={}Z+\langle X^{-1},Z\rangle
       (H_\alpha-2X-X^{-1}-H_\alpha+X)\\
  &\quad={}Z-\langle X^{-1},Z\rangle X
       -\langle X^{-1},Z\rangle X^{-1}.
\end{aligned}
\]
But
\[
  Z=\langle X^{-1},Z\rangle X,
  \qquad
  \langle X^{-1},Z\rangle X^{-1}=X^{-2}Z.
\]
Thus
\[
  a_\alpha\bigl(w_\alpha(X)\bigr)^{-1}=-X^{-2},
\]
and consequently
\[
  a_\alpha\bigl(w_\alpha(X)\bigr)=-X^2.
\]

\paragraph{\textbf{Corollary 3.7.}}
\SGARefTarget{sga3:xx:3.7}
\emph{In particular, one has}
\[
  \intaut\bigl(w_\alpha(X)\bigr)\exp(X)=\exp(-X^{-1}).
\]
\emph{Hence, by the definition of \(w_\alpha(X)\),}
\[
  w_\alpha(X)\exp(X)w_\alpha(X)^{-1}
    =\exp(-X)w_\alpha(X)\exp(-X),
\]
\emph{and therefore, by an immediate calculation,}
\[
  \bigl(w_\alpha(X)\exp(X)\bigr)^3=e.
\]

\paragraph{\textbf{Corollary 3.8.}}
\SGARefTarget{sga3:xx:3.8}
\emph{Let \(X\in\Gamma(S,\mathfrak g^\alpha)^\times\), and let
\(n\in\mathbb Z\), \(n\ne0\).  Then \(w_\alpha(X)\) is the unique
section \(w\in G(S)\) satisfying}

\smallskip
\noindent\emph{\textup{\SGAInlineRefTarget{sga3:xx:3.8:item:i}(i)}
\(\intaut(w)\circ n\alpha^\ast=-n\alpha^\ast\);}

\smallskip
\noindent\emph{\textup{\SGAInlineRefTarget{sga3:xx:3.8:item:ii}(ii)}
\(\bigl(w\exp(X)\bigr)^3=e\).}

\smallskip
We know that \(w_\alpha(X)\) satisfies these conditions.  Conversely,
let \(w\in G(S)\) satisfy \textup{\SGARefLink{sga3:xx:3.8:item:i}{(i)}} and \textup{\SGARefLink{sga3:xx:3.8:item:ii}{(ii)}}.  By \SGARefLink{sga3:xx:3.2}{3.2} and
\SGARefLink{sga3:xx:3.1}{3.1} \textup{\SGARefLink{sga3:xx:3.2:item:ii}{(ii)}}, there exists \(t\in T(S)\) such that
\[
  w=w_\alpha(X)t.
\]
Set \(u=\exp(X)\).  Then
\[
\begin{aligned}
  wuw^{-1}
    &=w_\alpha(X)t\exp(X)t^{-1}w_\alpha(X)^{-1}\\
    &=\exp\bigl(-\alpha(t)X^{-1}\bigr),
\end{aligned}
\]
whereas
\[
\begin{aligned}
  u^{-1}wu^{-1}
    &={}\exp(-X)w_\alpha(X)t\exp(-X)\\
    &={}\exp(-X)w_\alpha(X)\exp(-X)
       \exp\bigl(X-\alpha(t)X\bigr)t\\
    &={}\exp(-X^{-1})\exp\bigl(X-\alpha(t)X\bigr)t\\
    &={}\exp(-X^{-1})t
       \exp\bigl((\alpha(t)^{-1}-1)X\bigr).
\end{aligned}
\]
Now
\[
  (wu)^3=e
  \quad\Longleftrightarrow\quad
  wuw^{-1}=u^{-1}wu^{-1}.
\]
Comparing the two decompositions of this element in
\[
  U_{-\alpha}\cdot T\cdot U_\alpha
\]
gives \(t=e\).
